\documentclass[10pt,a4paper]{amsart}
\usepackage[margin=19mm]{geometry}
\usepackage{amsmath,amssymb,mathtools}
\usepackage[scr=rsfs,cal=euler]{mathalfa}
\usepackage{xfrac}
\usepackage[unicode,hidelinks,bookmarks=false]{hyperref}
\newcommand{\ie}{i.e.\ }
\newcommand{\cf}{cf.\ }
\newcommand{\resp}{respectively\ }
\newcommand{\Subsection}{\subsection}
\providecommand{\nobreakdash}{\nobreak\hbox{-}}
\setlength{\emergencystretch}{2em}
\multlinegap0pt
\usepackage[shortlabels]{enumitem}
\usepackage{xcolor}
\usepackage{amssymb}
\usepackage{mathtools}
\newtheorem{lemma}{Lemma}
\newtheorem{prop}{Proposition}
\def\Ee{\mathcal E}
\def\Le{\mathcal L}
\def\Me{\mathcal M}
\def\Ne{\mathcal N}
\def\Se{\mathcal S}
\title{Sufficient positive maps between von Neumann algebras: R{\'e}nyi divergences}
\author{Anna Jen\v{c}ov\'a, Lauritz van Luijk}
\date{Source: arXiv:2608.28510v1 (CC BY 4.0)}
\begin{document}
\maketitle

\noindent\emph{Source and licence.} Sufficient positive maps between von Neumann algebras: R{\'e}nyi divergences. Version arXiv:2608.28510v1 (CC BY 4.0), distributed by arXiv under the Creative Commons Attribution 4.0 International licence, http://creativecommons.org/licenses/by/4.0/, as recorded in the arXiv licence metadata for this exact version and shown on the abstract page https://arxiv.org/abs/2608.28510v1. Full licence notice: \texttt{licenses/arxiv-2608.28510-CC-BY-4.0.txt}.\par\smallskip

\noindent\textit{Editorial scope.} Counted unit: the article's prop prop:universal (source0.tex lines 476--483). The article's complete original proof of it is delivered with it (source0.tex lines 485--509, one \texttt{proof} environment of 25 lines). No mathematics is written, repaired or re-derived: the statement and the proof are the article's own lines, in the article's own notation, and no retained line is substituted or re-flowed.  Retained uncounted as the source-local context the proof needs: lines 363--368.  Nothing was omitted from any retained span, so the package carries no omission record.  No retained line contains a citation, so the wrapper carries no bibliography and no bibliography entry.  No retained line contains a figure environment, an image-inclusion command or drawing code, so no image is delivered and the package declares no asset dependency.  Editorial formatting only, in addition: an AMS \texttt{amsart} preamble that restates the article's own declarations and macros used by the retained lines, namely the environments \texttt{lemma}, \texttt{prop} on separate counters, together with the standard packages those lines need.  The retained environments print in this document as Lemma 1, Proposition 1 in order of appearance, whereas the article prints the same items with the numbering of its own full document.  The complete native source of the article as distributed in the arXiv e-print for this exact version is bundled unchanged as \texttt{sources/208794/source0.tex}.
\begin{lemma}\label{lemma:idempotent} Let $E:\Me\to \Me$ be a conditional expectation.
Let $\omega$ be a normal faithful state such that $\omega\circ E=\omega$. 
We then have $E_\omega=E$, and $E_2$ is an
 orthogonal projection on $L^2(\Me)$. 

\end{lemma}

\begin{prop}\label{prop:universal} 
An NUP map  $T:\Ne\to \Me$ is  sufficient with respect
to $\Ee$ if and only if
\[
\rho\circ T\circ T_\omega=\rho,\qquad \forall \rho\in \Se.
\]

\end{prop}

\begin{proof} Assume that $T$ is sufficient and let $S:\Me\to \Ne$ be a recovery map, then
clearly $T\circ S\in \Le$, and therefore
\[
T\circ S\circ E=E\circ T\circ S=E.
\]
Using the representation as contractions on $L^2(\Me)$ with the faithful state $\omega$, we obtain
\[
S_2^*T_2^*E_2=E_2,
\]
since $E_2$ is an orthogonal projection on $L^2(\Me)$, by Lemma \ref{lemma:idempotent}. 
Since both $T_2$ and $S_2$ are contractions, we have 
\[
\|\xi\|_2=\|S_2^*T_2^*\xi\|_2\le \|T_2^*\xi\|_2\le \|\xi\|_2,
\]
and hence  $\|T_2^*\xi\|_2=\|\xi\|_2$, for any $\xi$ in the range of $E_2$. This is
equivalent to $T_2T_2^*E_2=E_2$. 
Since $E=E_2^*$ and $T_2T_2^*$ are self-adjoint, we also have $E_2T_2T_2^*=E_2$, which shows $E\circ T\circ T_\omega=E$. 
But then for
any $\rho\in \Se$, we get
\[
\rho\circ T\circ T_\omega=\rho\circ E\circ T\circ T_\omega=\rho\circ E=\rho.
\]
The converse statement is clear.

\end{proof}

\end{document}
